% =====================================================================
%  diagrams/estructura-29119.tex
%  Figura: estructura general de ISO/IEC/IEEE 29119 y su relación
%          con las partes de este documento.
%
%  CORRECCIONES respecto del diagrama original («normal general.png»):
%   1. El original colgaba «Procesos organizacionales · política y
%      estrategia» de la Parte 1. Los procesos —incluido el
%      organizacional— pertenecen a la Parte 2; la Parte 1 aporta
%      conceptos generales y vocabulario.
%   2. Se elimina toda sugerencia de que una parte sustituye o contiene
%      a otra: son partes complementarias de una misma norma.
%   3. Se añade la correspondencia explícita con las secciones y anexos
%      de este documento, que el original no representaba.
%
%  Nota de mantenimiento: los nodos se colocan con «positioning»
%  (below=… of …) en lugar de coordenadas absolutas, de modo que al
%  cambiar el texto de un cuadro los demás se desplacen y nunca se
%  superpongan.
% =====================================================================
\begin{figure}[htbp]
\centering
\begin{tikzpicture}[
  font=\scriptsize,
  node distance=4mm,
  raiz/.style   ={draw=uvazul, very thick, fill=uvazul!12, rounded corners=3pt,
                  text width=5.0cm, align=center, inner sep=4pt},
  parte/.style  ={draw=uvazul, thick, fill=uvazul!7, rounded corners=2pt,
                  text width=2.35cm, align=center, inner sep=4pt},
  sub/.style    ={draw=gray!70, fill=gray!8, rounded corners=2pt,
                  text width=2.35cm, align=center, inner sep=4pt},
  doc/.style    ={draw=instrverde, thick, fill=instrverde!8, rounded corners=2pt,
                  text width=3.7cm, align=center, inner sep=4pt},
  fl/.style     ={-{Latex[length=4pt]}, draw=uvazul!80},
  mapa/.style   ={-{Latex[length=4pt]}, draw=instrverde!85, dashed}
]

% ------------------------------- raíz --------------------------------
\node[raiz] (n) {\textbf{ISO/IEC/IEEE 29119}\\ Prueba de software (norma multiparte)};

% ------------------------------ partes -------------------------------
\node[parte, below=10mm of n] (p3) {\textbf{Parte 3}\\ Documentación de prueba};
\node[parte, left=6mm of p3]  (p2) {\textbf{Parte 2}\\ Procesos de prueba};
\node[parte, left=6mm of p2]  (p1) {\textbf{Parte 1}\\ Conceptos generales y vocabulario};
\node[parte, right=6mm of p3] (p4) {\textbf{Parte 4}\\ Técnicas de prueba};
\node[parte, right=6mm of p4] (p5) {\textbf{Parte 5}\\ Pruebas por palabras clave};

\foreach \d in {p1,p2,p3,p4,p5} { \draw[fl] (n) -- (\d); }

% --------------------------- contenido -------------------------------
% Parte 1
\node[sub, below=of p1] (v1) {Vocabulario y conceptos que usan las demás partes};
\draw[fl] (p1) -- (v1);

% Parte 2: las tres capas de procesos
\node[sub, below=of p2] (c1) {Capa organizacional};
\node[sub, below=of c1] (c2) {Gestión de la prueba};
\node[sub, below=of c2] (c3) {Procesos dinámicos};
\draw[fl] (p2) -- (c1); \draw[fl] (c1) -- (c2); \draw[fl] (c2) -- (c3);

% Parte 3: tipos de documentación
\node[sub, below=of p3] (d1) {Documentos de gestión};
\node[sub, below=of d1] (d2) {Especificaciones de nivel};
\node[sub, below=of d2] (d3) {Registros de ejecución e incidentes};
\draw[fl] (p3) -- (d1); \draw[fl] (d1) -- (d2); \draw[fl] (d2) -- (d3);

% Parte 4: familias de técnicas
\node[sub, below=of p4] (t1) {Basadas en especificación};
\node[sub, below=of t1] (t2) {Basadas en estructura};
\node[sub, below=of t2] (t3) {Basadas en experiencia};
\draw[fl] (p4) -- (t1); \draw[fl] (t1) -- (t2); \draw[fl] (t2) -- (t3);

% Parte 5
\node[sub, below=of p5] (k1) {Enfoque opcional; se usa sólo si el proyecto lo adopta};
\draw[fl] (p5) -- (k1);

% ------------------- correspondencia con este documento --------------
\node[fit=(v1)(c3)(d3)(t3)(k1), inner sep=0pt, draw=none] (todos) {};

\node[align=center, font=\scriptsize\bfseries, text=instrverde,
      below=7mm of todos.south] (rot) {ESTE DOCUMENTO};

\node[doc, below=3mm of rot] (e2) {\textbf{Parte II}\\ Reporte de finalización};
\node[doc, left=5mm of e2]   (e1) {\textbf{Parte I}\\ Plan de prueba (secciones 1--9)};
\node[doc, right=5mm of e2]  (e3) {\textbf{Anexos A--G}\\ casos, procedimientos, ejecuciones, incidentes, informes, entorno, trazabilidad};

\draw[mapa] (c2.west) to[out=190,in=100] (e1.north);
\draw[mapa] (d1.east) to[out=0,in=90]    (e2.north east);
\draw[mapa] (d3.east) to[out=350,in=110] (e3.north west);
\draw[mapa] (t3.south) to[out=280,in=80] (e3.north);

% ------------------------------ leyenda ------------------------------
\node[align=left, font=\scriptsize, text width=15.2cm, below=5mm of e2] (leg)
  {\textbf{Cómo leer la figura.} Las flechas azules indican \emph{qué contiene} cada parte de la
   norma; las verdes discontinuas, \emph{dónde} se refleja ese contenido en este documento.
   Ninguna parte sustituye a otra: el plan no reemplaza a los registros que producen los procesos
   dinámicos, ni las técnicas de la Parte~4 son obligatorias en su totalidad.};

\end{tikzpicture}
\caption[Estructura de ISO/IEC/IEEE 29119 y su relación con este documento]%
{Estructura general de \norma\ y correspondencia de sus partes con las secciones y anexos de
este documento.}
\label{fig:estructura}
\end{figure}

\begin{observacion}[Diagrama de estructura de la norma]
\obsproblema{El diagrama de referencia colgaba «Procesos organizacionales · política y
estrategia» de la Parte~1, y sugería una jerarquía en la que unas partes sustituían a otras.}
\obscambio{Los procesos —incluido el organizacional— cuelgan de la Parte~2; la Parte~1 aporta
conceptos y vocabulario. Se añade la correspondencia con las secciones y anexos de este
documento y una leyenda que aclara que ninguna parte reemplaza a otra.}
\obsfuente{OBS §1 · DIAG}
\end{observacion}
